% El modelo de Mundell (1961): perturbación asimétrica y permanente de demanda
% en una unión monetaria (España y Alemania) y los 2 mecanismos de ajuste.
% Mismos valores, estilos y colores que da-oa-abierta del 3.A.32 (y da-oa del
% 3.A.5): DA en morado, OA de corto plazo en azul, OA de largo plazo en verde
% y lo desplazado, discontinuo (la DA', en rojo).
%   Y^n = 7; DA: Y = 11/3 + 14/(P + 2); OA_CP(W_0): P = Y - 4,8; A = (7; 2,2)
%   a) Perturbación: DA'_ES = DA - 2, B_ES = (6,018; 1,218) sobre OA_CP(W_0)
%      ((P + 2)^2 + 1,133(P + 2) = 14); DA'_DE = DA + 4/3, B_DE = (7,8; 3)
%   b) Flexibilidad salarial: OA_CP(W_1) española P = Y - 6,375, C_ES =
%      (7; 0,625) (da(0,625) - 2 = 7); alemana P = Y - 2, C_DE = (7; 5)
%   c) Movilidad del trabajo: Y^n' = 5 en España, OA'_CP: P = Y - 2,8,
%      C_ES = (5; 2,2) (da(2,2) - 2 = 5); Y^n' = 8,333 en Alemania,
%      OA'_CP: P = Y - 6,133, C_DE = (8,333; 2,2) (da(2,2) + 4/3)
% Autor: Víctor Gutiérrez Marcos
\documentclass[tikz,border=4pt]{standalone}
% ==========================================================================
% tikz-tcee.tex — Estilos comunes de los gráficos TikZ del temario
% Autor: Víctor Gutiérrez Marcos
%
% Cada gráfico es un fichero standalone en fuentes/<tema>/graficos/:
%
%   \documentclass[tikz,border=4pt]{standalone}
%   % ==========================================================================
% tikz-tcee.tex — Estilos comunes de los gráficos TikZ del temario
% Autor: Víctor Gutiérrez Marcos
%
% Cada gráfico es un fichero standalone en fuentes/<tema>/graficos/:
%
%   \documentclass[tikz,border=4pt]{standalone}
%   % ==========================================================================
% tikz-tcee.tex — Estilos comunes de los gráficos TikZ del temario
% Autor: Víctor Gutiérrez Marcos
%
% Cada gráfico es un fichero standalone en fuentes/<tema>/graficos/:
%
%   \documentclass[tikz,border=4pt]{standalone}
%   \input{../../../latex/tikz-tcee}
%   \begin{document}
%   \begin{tikzpicture}[tcee]
%     \ejes{Q}{P}{6}{5}
%     \draw[curva1] (0.5,4.5) -- (5,0.8) node[etiqueta,right] {$D$};
%     \capa{2}{\draw[curva2] ...;}   % en el vídeo aparece en el paso 2
%   \end{tikzpicture}
%   \end{document}
%
% Paleta: la de la web (morado, dorado, salvia) más tres colores de apoyo
% distinguibles en blanco y negro por el trazo.
% ==========================================================================
\usepackage[T1]{fontenc}
\usepackage{newpxtext,newpxmath}
\usepackage{xcolor}
\usepackage{tikz,pgfplots}
\pgfplotsset{compat=1.18}
\usetikzlibrary{arrows.meta,calc,positioning,patterns,decorations.pathreplacing,intersections,shapes.misc,backgrounds}

\definecolor{tceemorado}{HTML}{5F2987}
\definecolor{tceedorado}{HTML}{B8860B}
\definecolor{tceeverde}{HTML}{3C7D3A}
\definecolor{tceeazul}{HTML}{1F5F99}
\definecolor{tceerojo}{HTML}{B22222}
\definecolor{tceenaranja}{HTML}{E07B28}
\definecolor{tceegris}{HTML}{767171}
\definecolor{tceesalvia}{HTML}{E2EFD9}

% Capas para el vídeo: \capa{n}{…} dibuja su contenido solo a partir del paso n.
% En el PDF, la web y Word se ve todo; generar-video.py compila el gráfico una
% vez por paso con \def\tceepaso{k} para que cada elemento aparezca al narrarlo.
% Vectores con flecha encima, como en el texto (tcee.sty): \vect{x}, \vect{p}
\providecommand{\vect}[1]{\vec{#1}}
\providecommand{\tceepaso}{1000}
\newcommand{\capa}[2]{\ifnum#1>\tceepaso\relax\else#2\fi}
% Gráficos por pasos del mismo tamaño (propuesta del 3.A.29, 11/11/2026): con
% \def\tceefantasma{1} antes del \input (p. ej., x-paso1.tex:
% «\def\tceepaso{3}\def\tceefantasma{1}\input{x.tex}»), las capas posteriores
% al paso se dibujan invisibles y ocupan su sitio. Sin \tceefantasma (el vídeo),
% no cambia nada.
\ifdefined\tceefantasma
  \renewcommand{\capa}[2]{\ifnum#1>\tceepaso\relax\begin{scope}[opacity=0]#2\end{scope}\else#2\fi}
\fi

\tikzset{
  tcee/.style={line cap=round,line join=round,font=\small,>={Stealth[length=2.2mm]}},
  eje/.style={->,thick,black},
  curva1/.style={very thick,tceemorado},
  curva2/.style={very thick,tceeazul},
  curva3/.style={very thick,tceeverde},
  curva4/.style={very thick,tceedorado},
  curvanueva/.style={very thick,tceerojo},
  desplazada/.style={very thick,dashed},
  guia/.style={thin,dashed,tceegris},
  punto/.style={circle,fill=black,inner sep=1.3pt},
  etiqueta/.style={font=\small,inner sep=2pt},
  flecha/.style={->,thick,tceerojo},
  area/.style={fill=tceemorado,fill opacity=0.18,draw=none},
  area2/.style={fill=tceedorado,fill opacity=0.22,draw=none},
}

% \ejes{etiqueta x}{etiqueta y}{largo x}{alto y}
\newcommand{\ejes}[4]{%
  \draw[eje] (0,0) -- (#3,0) node[below left=2pt and 0pt] {#1};
  \draw[eje] (0,0) -- (0,#4) node[left] {#2};
  \node[below left] at (0,0) {$0$};}
% \proyeccion{(x,y)}{etq x}{etq y}: líneas guía a los ejes con etiquetas
\newcommand{\proyeccion}[3]{%
  \draw[guia] let \p1=#1 in (\x1,0) node[below] {#2} |- (0,\y1) node[left] {#3};}

% --------------------------------------------------------------------------
% Paleta «dictamen» (dic…): gráficos de datos de la Guía del dictamen (1D01) y
% de cualquier gráfico de coyuntura que quiera los mismos colores. Un concepto,
% un color en todos los gráficos (CRITERIOS §7, «tres reglas», 24/10/2026);
% tabla completa en fuentes/1D01/_trabajo/paleta-1D01.md. No cambia nada de
% lo anterior: son colores y estilos nuevos.
% Actividad: el PIB (y sus variantes: PIB por persona, PIB mundial), morado.
\definecolor{dicpib}{HTML}{5F2987}
% Demanda interna, gama cálida: consumo (naranjas: privado oscuro, público
% claro), inversión (verde) y existencias (gris con trama, aparte).
\definecolor{dicconspriv}{HTML}{C76A1F}
\definecolor{dicconspub}{HTML}{F5C690}
\definecolor{dicfbcf}{HTML}{2A6B3C}
\definecolor{dicexistencias}{HTML}{B4B4B4}
% Demanda externa y comercio, gama fría (azul claro, que se separa del morado)
\definecolor{dicexterna}{HTML}{6AA6DA}
% Zonas: zona del euro (azul marino), España (rojo), resto de países (grises)
\definecolor{dicze}{HTML}{17406D}
\definecolor{dicespana}{HTML}{C8382E}
\definecolor{dicotro}{HTML}{8A8A8A}
\definecolor{dicotroclaro}{HTML}{B5B5B5}
\definecolor{dicotromasclaro}{HTML}{DADADA}
% Precios: inflación general (negro, continuo) y subyacente (negro,
% discontinuo); componentes: energía (bermellón), alimentos (amarillo) y los
% de la subyacente en azules (bienes industriales claro, servicios oscuro).
\definecolor{dicinflacion}{HTML}{1A1A1A}
\definecolor{dicenergia}{HTML}{D55E00}
\definecolor{dicalimentos}{HTML}{F2D45C}
\definecolor{dicbine}{HTML}{7FB2E0}
\definecolor{dicservicios}{HTML}{1F4E89}
% Trabajo: empleo (verde azulado), horas (el mismo, claro), tasa de actividad
% (verde azulado oscuro), paro (teja); productividad (dorado de la web).
\definecolor{dicempleo}{HTML}{1B7F6B}
\definecolor{dichoras}{HTML}{86CBB8}
\definecolor{dicactividad}{HTML}{0F5247}
\definecolor{dicparo}{HTML}{A23B2A}
\definecolor{dicproductividad}{HTML}{B8860B}
% Hogares: renta disponible (pizarra) y tasa de ahorro (negro); el consumo es
% dicconspriv.
\definecolor{dicrenta}{HTML}{4A5A6A}
\definecolor{dicahorro}{HTML}{1A1A1A}
% Política monetaria: tipos del BCE en azules (depósito, el de la zona del
% euro; operaciones principales, medio; facilidad marginal, claro); balance:
% préstamos a los bancos (dorado claro), resto del activo (gris claro).
\definecolor{dicdeposito}{HTML}{17406D}
\definecolor{dicopf}{HTML}{4F86C6}
\definecolor{dicfmc}{HTML}{9EC5EA}
\definecolor{dicprestamos}{HTML}{D9A520}
\definecolor{dicresto}{HTML}{CFCFCF}
% Ciclo: series adelantadas, coincidentes (el PIB) y retrasadas
\definecolor{dicadelantada}{HTML}{D55E00}
\definecolor{dicretrasada}{HTML}{6AA6DA}
% Referencias: rejilla, previsión, sombreados
\definecolor{dicreferencia}{HTML}{767171}

\tikzset{
  % barra de un componente: \addplot[dicbarra=dicfbcf]
  dicbarra/.style={fill=#1, draw=#1!65!black, line width=0.3pt},
  % existencias: gris con trama, para que se distinga también en gris
  dicbarratrama/.style={fill=#1, draw=#1!65!black, line width=0.3pt,
    postaction={pattern=north east lines, pattern color=white}},
  % línea de una serie; con halo blanco cuando pasa sobre barras
  diclinea/.style={very thick, #1},
  diclineahalo/.style={very thick, #1, preaction={draw=white, line width=3.4pt, solid}},
  % puntos de una serie (capa aparte, después de la línea)
  dicpuntos/.style={only marks, mark=*, mark size=1.8pt,
    mark options={fill=white, draw=#1, line width=0.9pt, solid}},
  % línea del cero o de un umbral, después de las barras
  diccero/.style={thin, black},
  dicumbral/.style={thin, dashed, dicreferencia},
}
% Capas de los gráficos de datos (CRITERIOS §7, «tres reglas»), de abajo
% arriba: cuadrícula → sombreados (dicfondo) → barras y áreas (dicbarras) → ejes, cero y umbrales
% (axis lines) → líneas (main) → puntos (dicpuntos) → rótulos (axis
% descriptions). No basta el orden del código: pgfplots dibuja cuadrícula,
% ejes y marcas en sus propias capas (y area style pone los ejes encima), y
% dentro de un axis los \draw, \fill y \node no admiten capa (van con los
% datos de main); por eso todo lo que se dibuja dentro del axis es un \addplot.
% Uso: \begin{axis}[…, dicapas] (al final de las opciones);
%   barras y áreas: \addplot[on layer=dicbarras, …];
%   sombreados:     \dicsombra[estilo]{x1}{y1}{x2}{y2} (rectángulo, en dicfondo,
%                   bajo las barras y sobre la cuadrícula);
%   cero y umbrales: \dicraya[estilo]{(x1,y1) (x2,y2)} (en axis lines);
%   líneas: \addplot normal (main); puntos: \addplot[on layer=dicpuntos, …]
%   (las marcas de cualquier serie van solas a dicpuntos);
%   rótulos: \dicrotulo[opciones del nodo]{(x,y)}{texto} (en axis descriptions),
%   siempre sin fondo (nada de fill=white: si chocan con una línea, se desplazan;
%   CRITERIOS §7, regla 4).
% Los \capa{n} del vídeo no cambian: siguen envolviendo instrucciones completas.
\pgfdeclarelayer{dicfondo}
\pgfdeclarelayer{dicbarras}
\pgfdeclarelayer{dicpuntos}
\pgfplotsset{
  % (el segundo argumento reparte cuadrícula, ejes y rótulos como el estándar)
  layers/dic/.define layer set={axis background, axis grid, dicfondo, dicbarras,
    axis ticks, axis lines, axis tick labels, pre main, main, dicpuntos,
    axis descriptions, axis foreground}{/pgfplots/layers/standard},
  % leyenda sin caja blanca: toma el fondo del gráfico, también si el gráfico va
  % dentro de un cuadro gris (CRITERIOS §7, regla 4, 26/10/2026)
  dicapas/.style={axis on top=false, set layers=dic, mark layer=dicpuntos,
    legend style={fill=none}},
  dicsuelto/.style={forget plot, stack plots=false, sharp plot, mark=none},
}
\newcommand{\dicraya}[2][]{%
  \addplot[on layer=axis lines, dicsuelto, diccero, #1] coordinates {#2};}
\newcommand{\dicsombra}[5][]{%
  \addplot[on layer=dicfondo, dicsuelto, draw=none, #1]
    coordinates {(#2,#3) (#4,#3) (#4,#5) (#2,#5)};}
\newcommand{\dicrotulo}[3][]{%
  \addplot[on layer=axis descriptions, dicsuelto, draw=none] coordinates {#2}
    node[pos=0, #1] {#3};}

%   \begin{document}
%   \begin{tikzpicture}[tcee]
%     \ejes{Q}{P}{6}{5}
%     \draw[curva1] (0.5,4.5) -- (5,0.8) node[etiqueta,right] {$D$};
%     \capa{2}{\draw[curva2] ...;}   % en el vídeo aparece en el paso 2
%   \end{tikzpicture}
%   \end{document}
%
% Paleta: la de la web (morado, dorado, salvia) más tres colores de apoyo
% distinguibles en blanco y negro por el trazo.
% ==========================================================================
\usepackage[T1]{fontenc}
\usepackage{newpxtext,newpxmath}
\usepackage{xcolor}
\usepackage{tikz,pgfplots}
\pgfplotsset{compat=1.18}
\usetikzlibrary{arrows.meta,calc,positioning,patterns,decorations.pathreplacing,intersections,shapes.misc,backgrounds}

\definecolor{tceemorado}{HTML}{5F2987}
\definecolor{tceedorado}{HTML}{B8860B}
\definecolor{tceeverde}{HTML}{3C7D3A}
\definecolor{tceeazul}{HTML}{1F5F99}
\definecolor{tceerojo}{HTML}{B22222}
\definecolor{tceenaranja}{HTML}{E07B28}
\definecolor{tceegris}{HTML}{767171}
\definecolor{tceesalvia}{HTML}{E2EFD9}

% Capas para el vídeo: \capa{n}{…} dibuja su contenido solo a partir del paso n.
% En el PDF, la web y Word se ve todo; generar-video.py compila el gráfico una
% vez por paso con \def\tceepaso{k} para que cada elemento aparezca al narrarlo.
% Vectores con flecha encima, como en el texto (tcee.sty): \vect{x}, \vect{p}
\providecommand{\vect}[1]{\vec{#1}}
\providecommand{\tceepaso}{1000}
\newcommand{\capa}[2]{\ifnum#1>\tceepaso\relax\else#2\fi}
% Gráficos por pasos del mismo tamaño (propuesta del 3.A.29, 11/11/2026): con
% \def\tceefantasma{1} antes del \input (p. ej., x-paso1.tex:
% «\def\tceepaso{3}\def\tceefantasma{1}\input{x.tex}»), las capas posteriores
% al paso se dibujan invisibles y ocupan su sitio. Sin \tceefantasma (el vídeo),
% no cambia nada.
\ifdefined\tceefantasma
  \renewcommand{\capa}[2]{\ifnum#1>\tceepaso\relax\begin{scope}[opacity=0]#2\end{scope}\else#2\fi}
\fi

\tikzset{
  tcee/.style={line cap=round,line join=round,font=\small,>={Stealth[length=2.2mm]}},
  eje/.style={->,thick,black},
  curva1/.style={very thick,tceemorado},
  curva2/.style={very thick,tceeazul},
  curva3/.style={very thick,tceeverde},
  curva4/.style={very thick,tceedorado},
  curvanueva/.style={very thick,tceerojo},
  desplazada/.style={very thick,dashed},
  guia/.style={thin,dashed,tceegris},
  punto/.style={circle,fill=black,inner sep=1.3pt},
  etiqueta/.style={font=\small,inner sep=2pt},
  flecha/.style={->,thick,tceerojo},
  area/.style={fill=tceemorado,fill opacity=0.18,draw=none},
  area2/.style={fill=tceedorado,fill opacity=0.22,draw=none},
}

% \ejes{etiqueta x}{etiqueta y}{largo x}{alto y}
\newcommand{\ejes}[4]{%
  \draw[eje] (0,0) -- (#3,0) node[below left=2pt and 0pt] {#1};
  \draw[eje] (0,0) -- (0,#4) node[left] {#2};
  \node[below left] at (0,0) {$0$};}
% \proyeccion{(x,y)}{etq x}{etq y}: líneas guía a los ejes con etiquetas
\newcommand{\proyeccion}[3]{%
  \draw[guia] let \p1=#1 in (\x1,0) node[below] {#2} |- (0,\y1) node[left] {#3};}

% --------------------------------------------------------------------------
% Paleta «dictamen» (dic…): gráficos de datos de la Guía del dictamen (1D01) y
% de cualquier gráfico de coyuntura que quiera los mismos colores. Un concepto,
% un color en todos los gráficos (CRITERIOS §7, «tres reglas», 24/10/2026);
% tabla completa en fuentes/1D01/_trabajo/paleta-1D01.md. No cambia nada de
% lo anterior: son colores y estilos nuevos.
% Actividad: el PIB (y sus variantes: PIB por persona, PIB mundial), morado.
\definecolor{dicpib}{HTML}{5F2987}
% Demanda interna, gama cálida: consumo (naranjas: privado oscuro, público
% claro), inversión (verde) y existencias (gris con trama, aparte).
\definecolor{dicconspriv}{HTML}{C76A1F}
\definecolor{dicconspub}{HTML}{F5C690}
\definecolor{dicfbcf}{HTML}{2A6B3C}
\definecolor{dicexistencias}{HTML}{B4B4B4}
% Demanda externa y comercio, gama fría (azul claro, que se separa del morado)
\definecolor{dicexterna}{HTML}{6AA6DA}
% Zonas: zona del euro (azul marino), España (rojo), resto de países (grises)
\definecolor{dicze}{HTML}{17406D}
\definecolor{dicespana}{HTML}{C8382E}
\definecolor{dicotro}{HTML}{8A8A8A}
\definecolor{dicotroclaro}{HTML}{B5B5B5}
\definecolor{dicotromasclaro}{HTML}{DADADA}
% Precios: inflación general (negro, continuo) y subyacente (negro,
% discontinuo); componentes: energía (bermellón), alimentos (amarillo) y los
% de la subyacente en azules (bienes industriales claro, servicios oscuro).
\definecolor{dicinflacion}{HTML}{1A1A1A}
\definecolor{dicenergia}{HTML}{D55E00}
\definecolor{dicalimentos}{HTML}{F2D45C}
\definecolor{dicbine}{HTML}{7FB2E0}
\definecolor{dicservicios}{HTML}{1F4E89}
% Trabajo: empleo (verde azulado), horas (el mismo, claro), tasa de actividad
% (verde azulado oscuro), paro (teja); productividad (dorado de la web).
\definecolor{dicempleo}{HTML}{1B7F6B}
\definecolor{dichoras}{HTML}{86CBB8}
\definecolor{dicactividad}{HTML}{0F5247}
\definecolor{dicparo}{HTML}{A23B2A}
\definecolor{dicproductividad}{HTML}{B8860B}
% Hogares: renta disponible (pizarra) y tasa de ahorro (negro); el consumo es
% dicconspriv.
\definecolor{dicrenta}{HTML}{4A5A6A}
\definecolor{dicahorro}{HTML}{1A1A1A}
% Política monetaria: tipos del BCE en azules (depósito, el de la zona del
% euro; operaciones principales, medio; facilidad marginal, claro); balance:
% préstamos a los bancos (dorado claro), resto del activo (gris claro).
\definecolor{dicdeposito}{HTML}{17406D}
\definecolor{dicopf}{HTML}{4F86C6}
\definecolor{dicfmc}{HTML}{9EC5EA}
\definecolor{dicprestamos}{HTML}{D9A520}
\definecolor{dicresto}{HTML}{CFCFCF}
% Ciclo: series adelantadas, coincidentes (el PIB) y retrasadas
\definecolor{dicadelantada}{HTML}{D55E00}
\definecolor{dicretrasada}{HTML}{6AA6DA}
% Referencias: rejilla, previsión, sombreados
\definecolor{dicreferencia}{HTML}{767171}

\tikzset{
  % barra de un componente: \addplot[dicbarra=dicfbcf]
  dicbarra/.style={fill=#1, draw=#1!65!black, line width=0.3pt},
  % existencias: gris con trama, para que se distinga también en gris
  dicbarratrama/.style={fill=#1, draw=#1!65!black, line width=0.3pt,
    postaction={pattern=north east lines, pattern color=white}},
  % línea de una serie; con halo blanco cuando pasa sobre barras
  diclinea/.style={very thick, #1},
  diclineahalo/.style={very thick, #1, preaction={draw=white, line width=3.4pt, solid}},
  % puntos de una serie (capa aparte, después de la línea)
  dicpuntos/.style={only marks, mark=*, mark size=1.8pt,
    mark options={fill=white, draw=#1, line width=0.9pt, solid}},
  % línea del cero o de un umbral, después de las barras
  diccero/.style={thin, black},
  dicumbral/.style={thin, dashed, dicreferencia},
}
% Capas de los gráficos de datos (CRITERIOS §7, «tres reglas»), de abajo
% arriba: cuadrícula → sombreados (dicfondo) → barras y áreas (dicbarras) → ejes, cero y umbrales
% (axis lines) → líneas (main) → puntos (dicpuntos) → rótulos (axis
% descriptions). No basta el orden del código: pgfplots dibuja cuadrícula,
% ejes y marcas en sus propias capas (y area style pone los ejes encima), y
% dentro de un axis los \draw, \fill y \node no admiten capa (van con los
% datos de main); por eso todo lo que se dibuja dentro del axis es un \addplot.
% Uso: \begin{axis}[…, dicapas] (al final de las opciones);
%   barras y áreas: \addplot[on layer=dicbarras, …];
%   sombreados:     \dicsombra[estilo]{x1}{y1}{x2}{y2} (rectángulo, en dicfondo,
%                   bajo las barras y sobre la cuadrícula);
%   cero y umbrales: \dicraya[estilo]{(x1,y1) (x2,y2)} (en axis lines);
%   líneas: \addplot normal (main); puntos: \addplot[on layer=dicpuntos, …]
%   (las marcas de cualquier serie van solas a dicpuntos);
%   rótulos: \dicrotulo[opciones del nodo]{(x,y)}{texto} (en axis descriptions),
%   siempre sin fondo (nada de fill=white: si chocan con una línea, se desplazan;
%   CRITERIOS §7, regla 4).
% Los \capa{n} del vídeo no cambian: siguen envolviendo instrucciones completas.
\pgfdeclarelayer{dicfondo}
\pgfdeclarelayer{dicbarras}
\pgfdeclarelayer{dicpuntos}
\pgfplotsset{
  % (el segundo argumento reparte cuadrícula, ejes y rótulos como el estándar)
  layers/dic/.define layer set={axis background, axis grid, dicfondo, dicbarras,
    axis ticks, axis lines, axis tick labels, pre main, main, dicpuntos,
    axis descriptions, axis foreground}{/pgfplots/layers/standard},
  % leyenda sin caja blanca: toma el fondo del gráfico, también si el gráfico va
  % dentro de un cuadro gris (CRITERIOS §7, regla 4, 26/10/2026)
  dicapas/.style={axis on top=false, set layers=dic, mark layer=dicpuntos,
    legend style={fill=none}},
  dicsuelto/.style={forget plot, stack plots=false, sharp plot, mark=none},
}
\newcommand{\dicraya}[2][]{%
  \addplot[on layer=axis lines, dicsuelto, diccero, #1] coordinates {#2};}
\newcommand{\dicsombra}[5][]{%
  \addplot[on layer=dicfondo, dicsuelto, draw=none, #1]
    coordinates {(#2,#3) (#4,#3) (#4,#5) (#2,#5)};}
\newcommand{\dicrotulo}[3][]{%
  \addplot[on layer=axis descriptions, dicsuelto, draw=none] coordinates {#2}
    node[pos=0, #1] {#3};}

%   \begin{document}
%   \begin{tikzpicture}[tcee]
%     \ejes{Q}{P}{6}{5}
%     \draw[curva1] (0.5,4.5) -- (5,0.8) node[etiqueta,right] {$D$};
%     \capa{2}{\draw[curva2] ...;}   % en el vídeo aparece en el paso 2
%   \end{tikzpicture}
%   \end{document}
%
% Paleta: la de la web (morado, dorado, salvia) más tres colores de apoyo
% distinguibles en blanco y negro por el trazo.
% ==========================================================================
\usepackage[T1]{fontenc}
\usepackage{newpxtext,newpxmath}
\usepackage{xcolor}
\usepackage{tikz,pgfplots}
\pgfplotsset{compat=1.18}
\usetikzlibrary{arrows.meta,calc,positioning,patterns,decorations.pathreplacing,intersections,shapes.misc,backgrounds}

\definecolor{tceemorado}{HTML}{5F2987}
\definecolor{tceedorado}{HTML}{B8860B}
\definecolor{tceeverde}{HTML}{3C7D3A}
\definecolor{tceeazul}{HTML}{1F5F99}
\definecolor{tceerojo}{HTML}{B22222}
\definecolor{tceenaranja}{HTML}{E07B28}
\definecolor{tceegris}{HTML}{767171}
\definecolor{tceesalvia}{HTML}{E2EFD9}

% Capas para el vídeo: \capa{n}{…} dibuja su contenido solo a partir del paso n.
% En el PDF, la web y Word se ve todo; generar-video.py compila el gráfico una
% vez por paso con \def\tceepaso{k} para que cada elemento aparezca al narrarlo.
% Vectores con flecha encima, como en el texto (tcee.sty): \vect{x}, \vect{p}
\providecommand{\vect}[1]{\vec{#1}}
\providecommand{\tceepaso}{1000}
\newcommand{\capa}[2]{\ifnum#1>\tceepaso\relax\else#2\fi}
% Gráficos por pasos del mismo tamaño (propuesta del 3.A.29, 11/11/2026): con
% \def\tceefantasma{1} antes del \input (p. ej., x-paso1.tex:
% «\def\tceepaso{3}\def\tceefantasma{1}\input{x.tex}»), las capas posteriores
% al paso se dibujan invisibles y ocupan su sitio. Sin \tceefantasma (el vídeo),
% no cambia nada.
\ifdefined\tceefantasma
  \renewcommand{\capa}[2]{\ifnum#1>\tceepaso\relax\begin{scope}[opacity=0]#2\end{scope}\else#2\fi}
\fi

\tikzset{
  tcee/.style={line cap=round,line join=round,font=\small,>={Stealth[length=2.2mm]}},
  eje/.style={->,thick,black},
  curva1/.style={very thick,tceemorado},
  curva2/.style={very thick,tceeazul},
  curva3/.style={very thick,tceeverde},
  curva4/.style={very thick,tceedorado},
  curvanueva/.style={very thick,tceerojo},
  desplazada/.style={very thick,dashed},
  guia/.style={thin,dashed,tceegris},
  punto/.style={circle,fill=black,inner sep=1.3pt},
  etiqueta/.style={font=\small,inner sep=2pt},
  flecha/.style={->,thick,tceerojo},
  area/.style={fill=tceemorado,fill opacity=0.18,draw=none},
  area2/.style={fill=tceedorado,fill opacity=0.22,draw=none},
}

% \ejes{etiqueta x}{etiqueta y}{largo x}{alto y}
\newcommand{\ejes}[4]{%
  \draw[eje] (0,0) -- (#3,0) node[below left=2pt and 0pt] {#1};
  \draw[eje] (0,0) -- (0,#4) node[left] {#2};
  \node[below left] at (0,0) {$0$};}
% \proyeccion{(x,y)}{etq x}{etq y}: líneas guía a los ejes con etiquetas
\newcommand{\proyeccion}[3]{%
  \draw[guia] let \p1=#1 in (\x1,0) node[below] {#2} |- (0,\y1) node[left] {#3};}

% --------------------------------------------------------------------------
% Paleta «dictamen» (dic…): gráficos de datos de la Guía del dictamen (1D01) y
% de cualquier gráfico de coyuntura que quiera los mismos colores. Un concepto,
% un color en todos los gráficos (CRITERIOS §7, «tres reglas», 24/10/2026);
% tabla completa en fuentes/1D01/_trabajo/paleta-1D01.md. No cambia nada de
% lo anterior: son colores y estilos nuevos.
% Actividad: el PIB (y sus variantes: PIB por persona, PIB mundial), morado.
\definecolor{dicpib}{HTML}{5F2987}
% Demanda interna, gama cálida: consumo (naranjas: privado oscuro, público
% claro), inversión (verde) y existencias (gris con trama, aparte).
\definecolor{dicconspriv}{HTML}{C76A1F}
\definecolor{dicconspub}{HTML}{F5C690}
\definecolor{dicfbcf}{HTML}{2A6B3C}
\definecolor{dicexistencias}{HTML}{B4B4B4}
% Demanda externa y comercio, gama fría (azul claro, que se separa del morado)
\definecolor{dicexterna}{HTML}{6AA6DA}
% Zonas: zona del euro (azul marino), España (rojo), resto de países (grises)
\definecolor{dicze}{HTML}{17406D}
\definecolor{dicespana}{HTML}{C8382E}
\definecolor{dicotro}{HTML}{8A8A8A}
\definecolor{dicotroclaro}{HTML}{B5B5B5}
\definecolor{dicotromasclaro}{HTML}{DADADA}
% Precios: inflación general (negro, continuo) y subyacente (negro,
% discontinuo); componentes: energía (bermellón), alimentos (amarillo) y los
% de la subyacente en azules (bienes industriales claro, servicios oscuro).
\definecolor{dicinflacion}{HTML}{1A1A1A}
\definecolor{dicenergia}{HTML}{D55E00}
\definecolor{dicalimentos}{HTML}{F2D45C}
\definecolor{dicbine}{HTML}{7FB2E0}
\definecolor{dicservicios}{HTML}{1F4E89}
% Trabajo: empleo (verde azulado), horas (el mismo, claro), tasa de actividad
% (verde azulado oscuro), paro (teja); productividad (dorado de la web).
\definecolor{dicempleo}{HTML}{1B7F6B}
\definecolor{dichoras}{HTML}{86CBB8}
\definecolor{dicactividad}{HTML}{0F5247}
\definecolor{dicparo}{HTML}{A23B2A}
\definecolor{dicproductividad}{HTML}{B8860B}
% Hogares: renta disponible (pizarra) y tasa de ahorro (negro); el consumo es
% dicconspriv.
\definecolor{dicrenta}{HTML}{4A5A6A}
\definecolor{dicahorro}{HTML}{1A1A1A}
% Política monetaria: tipos del BCE en azules (depósito, el de la zona del
% euro; operaciones principales, medio; facilidad marginal, claro); balance:
% préstamos a los bancos (dorado claro), resto del activo (gris claro).
\definecolor{dicdeposito}{HTML}{17406D}
\definecolor{dicopf}{HTML}{4F86C6}
\definecolor{dicfmc}{HTML}{9EC5EA}
\definecolor{dicprestamos}{HTML}{D9A520}
\definecolor{dicresto}{HTML}{CFCFCF}
% Ciclo: series adelantadas, coincidentes (el PIB) y retrasadas
\definecolor{dicadelantada}{HTML}{D55E00}
\definecolor{dicretrasada}{HTML}{6AA6DA}
% Referencias: rejilla, previsión, sombreados
\definecolor{dicreferencia}{HTML}{767171}

\tikzset{
  % barra de un componente: \addplot[dicbarra=dicfbcf]
  dicbarra/.style={fill=#1, draw=#1!65!black, line width=0.3pt},
  % existencias: gris con trama, para que se distinga también en gris
  dicbarratrama/.style={fill=#1, draw=#1!65!black, line width=0.3pt,
    postaction={pattern=north east lines, pattern color=white}},
  % línea de una serie; con halo blanco cuando pasa sobre barras
  diclinea/.style={very thick, #1},
  diclineahalo/.style={very thick, #1, preaction={draw=white, line width=3.4pt, solid}},
  % puntos de una serie (capa aparte, después de la línea)
  dicpuntos/.style={only marks, mark=*, mark size=1.8pt,
    mark options={fill=white, draw=#1, line width=0.9pt, solid}},
  % línea del cero o de un umbral, después de las barras
  diccero/.style={thin, black},
  dicumbral/.style={thin, dashed, dicreferencia},
}
% Capas de los gráficos de datos (CRITERIOS §7, «tres reglas»), de abajo
% arriba: cuadrícula → sombreados (dicfondo) → barras y áreas (dicbarras) → ejes, cero y umbrales
% (axis lines) → líneas (main) → puntos (dicpuntos) → rótulos (axis
% descriptions). No basta el orden del código: pgfplots dibuja cuadrícula,
% ejes y marcas en sus propias capas (y area style pone los ejes encima), y
% dentro de un axis los \draw, \fill y \node no admiten capa (van con los
% datos de main); por eso todo lo que se dibuja dentro del axis es un \addplot.
% Uso: \begin{axis}[…, dicapas] (al final de las opciones);
%   barras y áreas: \addplot[on layer=dicbarras, …];
%   sombreados:     \dicsombra[estilo]{x1}{y1}{x2}{y2} (rectángulo, en dicfondo,
%                   bajo las barras y sobre la cuadrícula);
%   cero y umbrales: \dicraya[estilo]{(x1,y1) (x2,y2)} (en axis lines);
%   líneas: \addplot normal (main); puntos: \addplot[on layer=dicpuntos, …]
%   (las marcas de cualquier serie van solas a dicpuntos);
%   rótulos: \dicrotulo[opciones del nodo]{(x,y)}{texto} (en axis descriptions),
%   siempre sin fondo (nada de fill=white: si chocan con una línea, se desplazan;
%   CRITERIOS §7, regla 4).
% Los \capa{n} del vídeo no cambian: siguen envolviendo instrucciones completas.
\pgfdeclarelayer{dicfondo}
\pgfdeclarelayer{dicbarras}
\pgfdeclarelayer{dicpuntos}
\pgfplotsset{
  % (el segundo argumento reparte cuadrícula, ejes y rótulos como el estándar)
  layers/dic/.define layer set={axis background, axis grid, dicfondo, dicbarras,
    axis ticks, axis lines, axis tick labels, pre main, main, dicpuntos,
    axis descriptions, axis foreground}{/pgfplots/layers/standard},
  % leyenda sin caja blanca: toma el fondo del gráfico, también si el gráfico va
  % dentro de un cuadro gris (CRITERIOS §7, regla 4, 26/10/2026)
  dicapas/.style={axis on top=false, set layers=dic, mark layer=dicpuntos,
    legend style={fill=none}},
  dicsuelto/.style={forget plot, stack plots=false, sharp plot, mark=none},
}
\newcommand{\dicraya}[2][]{%
  \addplot[on layer=axis lines, dicsuelto, diccero, #1] coordinates {#2};}
\newcommand{\dicsombra}[5][]{%
  \addplot[on layer=dicfondo, dicsuelto, draw=none, #1]
    coordinates {(#2,#3) (#4,#3) (#4,#5) (#2,#5)};}
\newcommand{\dicrotulo}[3][]{%
  \addplot[on layer=axis descriptions, dicsuelto, draw=none] coordinates {#2}
    node[pos=0, #1] {#3};}

\begin{document}
\begin{tikzpicture}[tcee,x=0.55cm,y=0.55cm,
    declare function={da(\p)=11/3+14/(\p+2);}]
  % Marco fijo para que todas las capas del vídeo tengan el mismo tamaño
  \path (-2.2,-26.4) rectangle (33.2,11.4);
  % ==================================================== fila a: perturbación
  % 1: ejes de los 2 paneles y títulos
  \capa{1}{\node[font=\bfseries] at (6.3,10.8) {España};}
  \capa{1}{\node[anchor=base west, font=\small\bfseries] at (-1.2,9.3) {a) Perturbación asimétrica de demanda};}
  \capa{1}{\ejes{$Y$}{$P$}{12.6}{7.6}}
  \begin{scope}[xshift=9.6cm]
  \capa{1}{\node[font=\bfseries] at (6.3,10.8) {Alemania};}
  \capa{1}{\ejes{$Y$}{$P$}{12.6}{7.6}}
  \end{scope}
  % 2: la OA de largo plazo, vertical en Y^n, en los 2 países
  \capa{2}{\draw[curva3] (7,0) node[below, text=black] {$Y^n$} -- (7,7.6) node[etiqueta, above, text=tceeverde] {OA$_{\mathrm{LP}}$};}
  \begin{scope}[xshift=9.6cm]
  \capa{2}{\draw[curva3] (7,0) node[below, text=black] {$Y^n$} -- (7,7.6) node[etiqueta, above, text=tceeverde] {OA$_{\mathrm{LP}}$};}
  \end{scope}
  % 3: DA y OA de corto plazo con W_0 en los 2 países; equilibrio inicial A
  \capa{3}{\draw[curva1] plot[domain=6.2:1.2, samples=60, smooth] ({da(\x)},\x)
      node[etiqueta, right, text=tceemorado] {DA$_{\mathrm{ES}}$};}
  \capa{3}{\draw[curva2] (5.3,0.5) -- (10.2,5.4) node[etiqueta, right, text=tceeazul] {OA$_{\mathrm{CP}}(W_0)$};}
  \capa{3}{\draw[guia] (7,2.2) -- (0,2.2) node[left, text=black] {$P_0$};}
  \capa{3}{\node[punto] at (7,2.2) {};}
  \capa{3}{\node[etiqueta, right=4pt] at (7,2.2) {$A$};}
  \begin{scope}[xshift=9.6cm]
  \capa{3}{\draw[curva1] plot[domain=1.2:6.2, samples=60, smooth] ({da(\x)},\x)
      node[etiqueta, above, text=tceemorado] {DA$_{\mathrm{DE}}$};}
  \capa{3}{\draw[curva2] (5.3,0.5) -- (10.2,5.4) node[etiqueta, right, text=tceeazul] {OA$_{\mathrm{CP}}(W_0)$};}
  \capa{3}{\draw[guia] (7,2.2) -- (0,2.2) node[left, text=black] {$P_0$};}
  \capa{3}{\node[punto] at (7,2.2) {};}
  \capa{3}{\node[etiqueta] at (6.62,1.3) {$A$};}
  \end{scope}
  % 4: España: la demanda cae (DA') y la economía pasa a B, con paro
  \capa{4}{\draw[curvanueva, desplazada] plot[domain=5.6:0.4, samples=60, smooth] ({da(\x)-2},\x)
      node[etiqueta, right, yshift=2pt, text=tceerojo] {DA$'_{\mathrm{ES}}$};}
  \capa{4}{\draw[guia] (6.018,1.218) -- (0,1.218) node[left, text=black] {$P_1$};}
  \capa{4}{\draw[guia] (6.018,0) node[below, text=black] {$Y_1$} -- (6.018,1.218);}
  \capa{4}{\node[punto] at (6.018,1.218) {};}
  \capa{4}{\node[etiqueta, left=2pt] at (6.018,1.218) {$B$};}
  \capa{4}{\node[etiqueta, anchor=north west, font=\footnotesize] at (0.3,7.6) {recesión: paro};}
  % 5: Alemania: la demanda sube (DA') y la economía pasa a B, con
  %    presiones inflacionistas
  \begin{scope}[xshift=9.6cm]
  \capa{5}{\draw[curvanueva, desplazada] plot[domain=6.3:1.0, samples=60, smooth] ({da(\x)+4/3},\x)
      node[etiqueta, right, text=tceerojo] {DA$'_{\mathrm{DE}}$};}
  \capa{5}{\draw[guia] (7.8,3) -- (0,3) node[left, text=black] {$P_1$};}
  \capa{5}{\draw[guia] (7.8,0) node[below, text=black] {$Y_1$} -- (7.8,3);}
  \capa{5}{\node[punto] at (7.8,3) {};}
  \capa{5}{\node[etiqueta, right=3pt] at (7.8,3) {$B$};}
  \capa{5}{\node[etiqueta, anchor=north west, align=left, font=\footnotesize] at (0.3,7.6) {presiones\\inflacionistas};}
  \end{scope}
  % ==================================================== fila b: flexibilidad salarial
  \begin{scope}[yshift=-6.9cm]
  % 6: ejes y curvas de partida (OA_LP, DA', OA_CP(W_0), B) en los 2 países
  \capa{6}{\node[anchor=base west, font=\small\bfseries] at (-1.2,9.3) {b) Flexibilidad salarial (devaluación interna)};}
  \capa{6}{\ejes{$Y$}{$P$}{12.6}{7.6}}
  \capa{6}{\draw[curva3] (7,0) node[below, text=black] {$Y^n$} -- (7,7.6) node[etiqueta, above, text=tceeverde] {OA$_{\mathrm{LP}}$};}
  \capa{6}{\draw[curvanueva, desplazada] plot[domain=5.6:0.4, samples=60, smooth] ({da(\x)-2},\x)
      node[etiqueta, right, yshift=2pt, text=tceerojo] {DA$'_{\mathrm{ES}}$};}
  \capa{6}{\draw[curva2] (5.3,0.5) -- (10.2,5.4) node[etiqueta, right, text=tceeazul] {OA$_{\mathrm{CP}}(W_0)$};}
  \capa{6}{\draw[guia] (6.018,1.218) -- (0,1.218) node[left, text=black] {$P_1$};}
  \capa{6}{\node[punto] at (6.018,1.218) {};}
  \capa{6}{\node[etiqueta, left=2pt] at (6.018,1.218) {$B$};}
  \begin{scope}[xshift=9.6cm]
  \capa{6}{\ejes{$Y$}{$P$}{12.6}{7.6}}
  \capa{6}{\draw[curva3] (7,0) node[below, text=black] {$Y^n$} -- (7,7.6) node[etiqueta, above, text=tceeverde] {OA$_{\mathrm{LP}}$};}
  \capa{6}{\draw[curvanueva, desplazada] plot[domain=6.3:1.0, samples=60, smooth] ({da(\x)+4/3},\x)
      node[etiqueta, right, text=tceerojo] {DA$'_{\mathrm{DE}}$};}
  \capa{6}{\draw[curva2] (5.3,0.5) -- (10.2,5.4) node[etiqueta, right, text=tceeazul] {OA$_{\mathrm{CP}}(W_0)$};}
  \capa{6}{\draw[guia] (7.8,3) -- (0,3) node[left, text=black] {$P_1$};}
  \capa{6}{\node[punto] at (7.8,3) {};}
  \capa{6}{\node[etiqueta, right=3pt] at (7.8,3) {$B$};}
  \end{scope}
  % 7: España: los salarios bajan, la OA de corto plazo baja (W_1 < W_0) y la
  %    economía vuelve a Y^n con un precio menor (C)
  \capa{7}{\draw[flecha] (9.3,4.5) -- (10.35,3.45);}
  \capa{7}{\draw[curva2, desplazada] (6.6,0.225) -- (10.6,4.225) node[etiqueta, right, text=tceeazul] {OA$_{\mathrm{CP}}(W_1)$};}
  \capa{7}{\draw[guia] (7,0.625) -- (0,0.625) node[left, text=black] {$P_2$};}
  \capa{7}{\node[punto] (bC) at (7,0.625) {};}
  \capa{7}{\node[etiqueta] at (6.5,0.3) {$C$};}
  \capa{7}{\node[etiqueta, anchor=north west, align=left, font=\footnotesize] at (0.3,7.6) {$W\downarrow$, $P\downarrow$, $q\uparrow$};}
  % 8: Alemania: los salarios suben, la OA de corto plazo sube y la economía
  %    vuelve a Y^n con un precio mayor (C)
  \begin{scope}[xshift=9.6cm]
  \capa{8}{\draw[flecha] (7.1,3.2) -- (6.0,4.3);}
  \capa{8}{\draw[curva2, desplazada] (2.5,0.5) -- (8.4,6.4) node[etiqueta, right, text=tceeazul] {OA$_{\mathrm{CP}}(W_1)$};}
  \capa{8}{\draw[guia] (7,5) -- (0,5) node[left, text=black] {$P_2$};}
  \capa{8}{\node[punto] (bC2) at (7,5) {};}
  \capa{8}{\node[etiqueta, above left=1pt and 1pt] at (bC2) {$C$};}
  \capa{8}{\node[etiqueta, anchor=north west, align=left, font=\footnotesize] at (0.3,7.6) {$W\uparrow$, $P\uparrow$, $q\uparrow$};}
  \end{scope}
  \end{scope}
  % ==================================================== fila c: movilidad del trabajo
  \begin{scope}[yshift=-13.8cm]
  % 9: ejes y curvas de partida (OA_LP, DA', OA_CP, B) en los 2 países
  \capa{9}{\node[anchor=base west, font=\small\bfseries] at (-1.2,9.3) {c) Movilidad del trabajo};}
  \capa{9}{\ejes{$Y$}{$P$}{12.6}{7.6}}
  \capa{9}{\draw[curva3] (7,0) node[below, text=black] {$Y^n$} -- (7,7.6) node[etiqueta, above right=0pt and -6pt, text=tceeverde] {OA$_{\mathrm{LP}}$};}
  \capa{9}{\draw[curvanueva, desplazada] plot[domain=5.6:0.4, samples=60, smooth] ({da(\x)-2},\x)
      node[etiqueta, right, yshift=2pt, text=tceerojo] {DA$'_{\mathrm{ES}}$};}
  \capa{9}{\draw[curva2] (5.3,0.5) -- (10.2,5.4) node[etiqueta, right, text=tceeazul] {OA$_{\mathrm{CP}}(W_0)$};}
  \capa{9}{\node[punto] at (6.018,1.218) {};}
  \capa{9}{\node[etiqueta, left=2pt] at (6.018,1.218) {$B$};}
  \begin{scope}[xshift=9.6cm]
  \capa{9}{\ejes{$Y$}{$P$}{12.6}{7.6}}
  \capa{9}{\draw[curva3] (7,0) node[below, text=black] {$Y^n$} -- (7,7.6) node[etiqueta, above left=0pt and -6pt, text=tceeverde] {OA$_{\mathrm{LP}}$};}
  \capa{9}{\draw[curvanueva, desplazada] plot[domain=6.3:1.0, samples=60, smooth] ({da(\x)+4/3},\x)
      node[etiqueta, right, text=tceerojo] {DA$'_{\mathrm{DE}}$};}
  \capa{9}{\draw[curva2] (5.3,0.5) -- (10.2,5.4) node[etiqueta, right, text=tceeazul] {OA$_{\mathrm{CP}}(W_0)$};}
  \capa{9}{\node[punto] at (7.8,3) {};}
  \capa{9}{\node[etiqueta, left=1pt] at (7.8,3) {$B$};}
  \end{scope}
  % 10: España: salen trabajadores, la OA de largo plazo y la de corto plazo
  %     se desplazan a la izquierda (Y^n') y la economía vuelve a P_0 (C)
  \capa{10}{\draw[flecha] (6.8,6.6) -- (5.2,6.6);}
  \capa{10}{\draw[curva3, desplazada] (5,0) node[below, text=black] {$Y^{n\prime}$} -- (5,7.6) node[etiqueta, above left=0pt and -6pt, text=tceeverde] {OA$'_{\mathrm{LP}}$};}
  \capa{10}{\draw[curva2, desplazada] (3.3,0.5) -- (9.2,6.4) node[etiqueta, right, text=tceeazul] {OA$'_{\mathrm{CP}}$};}
  \capa{10}{\draw[guia] (5,2.2) -- (0,2.2) node[left, text=black] {$P_0$};}
  \capa{10}{\node[punto] (cC) at (5,2.2) {};}
  \capa{10}{\node[etiqueta, left=3pt] at (cC) {$C$};}
  % 11: Alemania: llegan trabajadores, las 2 OA se desplazan a la derecha y la
  %     economía vuelve a P_0 sobre la DA' (C); flecha «emigración»
  \begin{scope}[xshift=9.6cm]
  \capa{11}{\draw[flecha] (7.2,6.6) -- (8.13,6.6);}
  \capa{11}{\draw[curva3, desplazada] (8.333,0) node[below, text=black] {$Y^{n\prime}$} -- (8.333,7.6) node[etiqueta, above right=0pt and -6pt, text=tceeverde] {OA$'_{\mathrm{LP}}$};}
  \capa{11}{\draw[curva2, desplazada] (6.4,0.267) -- (10.4,4.267) node[etiqueta, right, text=tceeazul] {OA$'_{\mathrm{CP}}$};}
  \capa{11}{\draw[guia] (8.333,2.2) -- (0,2.2) node[left, text=black] {$P_0$};}
  \capa{11}{\node[punto] (cC2) at (8.333,2.2) {};}
  \capa{11}{\node[etiqueta, right=3pt] at (cC2) {$C$};}
  \end{scope}
  \capa{11}{\draw[flecha, very thick] (12.4,3.6) -- (16.2,3.6) node[midway, above, font=\footnotesize, text=black] {emigración};}
  \end{scope}
\end{tikzpicture}
\end{document}
